\documentclass[conference]{IEEEtran}
\usepackage[utf8]{inputenc}
\usepackage[T1]{fontenc}
\usepackage{amsmath,amssymb}
\usepackage{booktabs}
\usepackage{multirow}
\usepackage{graphicx}
\usepackage[dvipsnames]{xcolor}
\usepackage{enumitem}
\usepackage[colorlinks=true,linkcolor=blue,citecolor=blue,urlcolor=blue]{hyperref}
\usepackage{cite}
\usepackage{balance}

% ---- draft machinery -------------------------------------------------------
% \ph{} : pending experimental result (fill after runs; each is bound to a
%         pre-registered hypothesis whose falsifier is stated in the text)
\newcommand{\ph}[1]{\textcolor{Gray}{\textbf{[#1]}}}
% \wj{} : author standpoint comment -- states the pre-registered position
%         (what we bet happens, and what each outcome would mean). Remove
%         before camera-ready by redefining to \relax.
\newcommand{\wj}[1]{\textcolor{RedOrange}{\footnotesize\textsf{[WJ: #1]}}}
% toggle for submission: uncomment to strip all \wj comments
% \renewcommand{\wj}[1]{}
% ----------------------------------------------------------------------------

\begin{document}

\title{Channel-Gated, Not Capability-Gated: Measuring Surplus Loss and Testing\\ Constraint-as-Protection in LLM Agent Bargaining}

\author{\IEEEauthorblockN{Anonymous Author(s)}
\IEEEauthorblockA{Anonymous Institution(s)\\
anonymous@example.org}}

\maketitle

\begin{abstract}
Large language model (LLM) agents are increasingly delegated bilateral negotiation on behalf of a principal, and a principal fielding a weaker or cheaper model may face a counterparty running a stronger one. How much bargaining surplus that weak agent loses as the capability gap widens---and whether a near-zero-cost defense exists---has not been measured: existing weak-vs-strong studies score adversarial interaction by win/loss outcomes or supervision quality, never by a continuous economic outcome against a computable reference split, and existing evidence that protocol-level restriction protects a principal always restricts the strong or untrusted side, never the weak party's own expressive channel. Our central claim, tested rather than assumed, is that surplus extraction is \emph{channel-gated}: how much of a capability gap converts into lost surplus depends on the action-space channel the bargaining protocol leaves open to the weaker party, not on the gap alone. We build a bilateral bargaining harness with closed-form Pareto frontiers and three computable reference splits, quantify capability gap two ways, and measure how a weak agent's normalized surplus loss scales with that gap. We then delete the weak agent's free-text channel, leaving only schema-validated structured offers, and test---via a same-model causal contrast, a bilateral both-constrained contrast that separates asymmetric shielding from a mere format effect, and a cross-tier substitution comparison---whether this restriction lowers surplus loss, whether the protection is asymmetric, and whether it can substitute for a tier of capability. In our experiments, structured offers reduce normalized surplus loss by \ph{XX.X\%} at matched capability gap, remove \ph{YY.Y\%} of the free-text gap-to-loss slope, and allow a lower-capability tier to match the next stronger unconstrained tier within a pre-registered non-inferiority margin of $0.05$ normalized-loss units. Every hypothesis is pre-registered with an explicit falsification condition and a pre-committed fallback framing. The actionable artifact is a deployment-time, configuration-level output-schema constraint: a protocol/interface design lever for multi-agent systems available at effectively zero marginal cost, independent of the counterparty's model.
\end{abstract}

\begin{IEEEkeywords}
LLM agents, multi-agent systems, negotiation, capability asymmetry, agent protocols, empirical software engineering, pre-registration
\end{IEEEkeywords}

\section{Introduction}
\label{sec:intro}

Autonomous LLM agents increasingly act as representatives: they bargain over prices, allocate shared resources, and reach agreements with other automated systems on behalf of a human or organizational principal. In such settings the agent is not merely producing text; the principal's economic outcome depends on the agreement it reaches. A natural deployment choice is a smaller, cheaper model---sufficient for routine negotiation and easy to run at scale. The difficulty is strategic: the principal controls only its own agent, not the counterparty's, so a weak principal-agent may face a stronger adversarial agent with better reasoning, persuasion, and offer selection. The asymmetry arises by accident whenever two parties pick different model tiers, and by design whenever one side is adversarial. It raises a basic, unresolved question: \emph{how much surplus does a weak LLM agent lose as the capability gap to its adversary grows, and is that loss inevitable?}

The intuitive answer is that the loss is fixed by the gap---a stronger model out-negotiates a weaker one, and the only remedy is a stronger model. This paper's central, tested-not-assumed claim is that this intuition is incomplete: \emph{surplus extraction is channel-gated, not only capability-gated}. How much of a capability gap converts into lost surplus is a function of the action-space channel the bargaining protocol leaves open to the weaker party, not of the gap alone. A stronger adversary must route its advantage through some action space---natural-language arguments, anchors, deadlines, threats, offer sequences---to turn it into surplus. Restricting the weak agent to schema-validated structured offers may deny part of that advantage any channel through which to convert into surplus, \emph{without altering either model's capability}. We measure the surplus a weak agent loses as a function of the gap, then delete the weak agent's free-text channel and test whether this near-zero-cost restriction lowers the loss, whether the protection is \emph{asymmetric} (shielding the weak side specifically) rather than a symmetric format effect, and whether it can stand in for a tier of capability.

We do not assume restriction helps. Deleting the weak agent's free-text field removes both its exposure to rhetorical exploitation \emph{and} its own legitimate defenses---explanation, commitment, strategic framing, stalling. Whether restriction protects the weak agent, harms it, or does nothing is not decidable from first principles. That undecidedness is precisely the value of the study: constraint-as-protection is a falsifiable deployment hypothesis, not a design slogan. If it holds, a near-zero-cost mitigation---constrain your own agent's output format---becomes available and its size is measurable; if it fails, capability itself is the operative lever and a budget-constrained principal is simply exposed.

Existing work supplies pieces of an answer but does not close the question. Scalable-oversight research provides the \emph{method}: weak-vs-strong adversarial outcomes scale predictably with a capability gap and can be fit from round-robin play \cite{ScalingLawsOversight}---but the outcome scored is an oversight win-rate, never a continuous economic quantity against a fair reference, so the technique transfers while the finding does not. AI-control research provides the \emph{phenomenon}: protocol-level restriction can protect a principal from a stronger or untrusted counterparty \cite{AIControl,ControlTax}---but the restriction always falls on the strong, untrusted side or the surrounding audit protocol, and the outcome is catastrophe probability, not a bargaining split; no control result restricts the \emph{weak party's own} channel as self-protection. Recent work on bilateral trade between LLMs provides the closest \emph{architecture}: a structured bargaining simulator that likewise separates binding offers from natural-language messages, whose 15{,}000-negotiation tournament finds that aggressive anchoring, calibrated concession, and temporal patience associate with the highest surplus share, and that weaker models capture surplus only where wide bargaining ranges mask their weakness \cite{BilateralTrade2604}---but these strategy--surplus associations are read off free play correlationally; nobody manipulates the weak side's channel as a treatment, scores loss against a computable reference split, or tests cross-tier substitution.

The sharpest precedent makes our question non-obvious rather than merely unstudied. Heddaya et al.\ run essentially our manipulation---alternating structured numeric offers versus unrestricted natural language---in a controlled human bargaining experiment, and find that the \emph{average agreed price is statistically identical} across the two conditions (their language treatment does move process metrics---agreement rate, speed, price variance---but not the price itself) \cite{LanguageOfBargaining2023}. Read naively, this predicts our intervention should do nothing. But their dyads are capability-\emph{symmetric}: there is no advantage for the channel to carry, so removing it changes nothing about who captures surplus. Our question is what happens to that null \emph{once a capability gap opens}. If the channel is the conduit through which a capability advantage becomes surplus, a manipulation inert at zero gap should turn protective precisely as the gap widens---a condition$\times$gap \emph{interaction}, not a main effect at gap${\approx}0$. We pre-register this as a binding gate rather than a soft check: our zero-gap self-play cells run under \emph{both} action-space conditions specifically so the condition-null is testable at $\Delta{\approx}0$; reproducing the human price-null there validates the interaction design, and failing to reproduce it---LLM dyads breaking the human format-null even at symmetric capability---is itself pre-committed as the primary finding, so no outcome strands the study.

Turning the question into a measurable one surfaces four obstacles, each of which we convert into a design commitment rather than an assumption. First, \emph{the gap is not a scalar}: capability is multi-dimensional, so any ``gap'' is a projection, and a general-benchmark delta and a negotiation-specific rating fit from play can order the same pairs differently \cite{GreatModelsThinkAlike}. We quantify the gap two ways and require every scaling claim to hold under both at once, with a pre-registered rank-concordance floor below which the construct is declared unstable and results are reported per-measure---so disagreement is a bounded, reported finding, not a silent escape hatch. Second, \emph{``lost surplus'' needs a fair reference}: we score every game instance against three independently computable reference splits (Nash, Kalai--Smorodinsky, equal split) and validate the metric against a zero-gap self-play null before trusting any nonzero reading---an existence check no prior weak-vs-strong study needs, because none scores a continuous surplus. Third, \emph{the net effect of restriction is not decidable a priori, and could be a mere format effect}: any protection observed for the weak side could be a symmetric benefit of structured offers rather than asymmetric shielding. We therefore test a same-model contrast, a cross-tier substitution comparison, \emph{and} a confirmatory bilateral both-constrained contrast whose job is to separate asymmetric self-shielding from a format effect, with a pre-committed frame switch if the two are indistinguishable. Fourth, \emph{two channels are confounded}: a stronger model out-negotiates on substance and out-persuades on rhetoric in the same turn; the both-constrained control plus tactic-tagged transcripts separate how much of any effect is rhetorical from how much is substantive. We further pre-register a process-metric panel (rounds-to-agreement, impasse rate, offer count) with a pinned estimand precedence---the primary loss estimand \emph{includes} impasses, since a protection that works by killing deals is not protection---so a ``protection'' reading cannot be a compositional artifact of which episodes reached agreement.

For software engineering readers, the actionable output of this study is not a result in bargaining economics but a \emph{protocol/interface design decision for multi-agent systems}: a deployable output-schema constraint on one's own agent, whose protective value---if it holds---is an engineering lever available at configuration time, at effectively zero marginal cost, and independent of the counterparty's model. That is the artifact a practitioner deploys, and the study is designed so that every deployment-relevant claim about it carries a pre-registered falsifier.

In summary, we contribute:
\begin{itemize}[leftmargin=1.4em]
\item \textbf{Channel-gated surplus extraction, formalized and tested}: the claim that how much of a capability gap becomes lost surplus is a function of the action-space channel the protocol leaves the weaker party---stated as a falsifiable hypothesis and tested via pre-registered contrasts, with the capability-symmetric human price-null \cite{LanguageOfBargaining2023} as its designed, binding departure baseline and the correlational tournament findings of \cite{BilateralTrade2604} as the phenomenon it causally decomposes.
\item \textbf{A capability-gap-controlled bargaining harness with a validated surplus metric}: a dual gap measure with a rank-concordance stability gate, three independently computable reference splits, and a zero-gap existence check that validates the surplus-loss metric before any causal claim is built on it.
\item \textbf{A minimal, single-field weak-side self-restriction, tested for protection and asymmetry}: deleting exactly the weak agent's free-text field, evaluated via a same-model causal contrast, a cross-tier substitution comparison, and a confirmatory both-constrained contrast that separates asymmetric self-shielding from a symmetric format effect, with a pre-committed frame switch between the two.
\item \textbf{A channel-attribution analysis} separating rhetorical exploitation from substantive bargaining skill, keeping a causal between-condition bound and a correlational tactic-level reading at their true, different evidence strengths.
\end{itemize}

\section{Problem Definition}
\label{sec:problem}

\subsection{Formal Setup}

Consider a directed model pair $(w,s)$, where $w$ is a weak-capability agent representing a principal and $s$ is a strong-capability adversary representing the counterparty. The pair negotiates over bargaining game instances $g$, each a multi-issue alternating-offers game with private integer valuations, reservation values $r_w$ and $r_s$ known to each side but not its counterpart, and a deadline of $T{=}10$ rounds. Each episode $e$ under instance $g$ and action-space condition $c \in \{\textsf{structured}, \textsf{free}\}$---where $c$ is applied to $w$'s turn only, except in a bilateral control where it is applied to both sides---yields an agreement flag and, if agreement is reached, a realized surplus split.

\textbf{Normalized surplus loss.} For the weak side,
\begin{equation}
L \;=\; \frac{u_{\mathrm{ref}} - u_{\mathrm{realized}}}{u_{\mathrm{ref}} - r_w},
\label{eq:loss}
\end{equation}
with impasse scored at $u_{\mathrm{realized}} = r_w$, hence $L{=}1$. The reference utility $u_{\mathrm{ref}}$ is the weak side's utility at the Nash-bargaining solution of $g$ (primary); the equal-split and Kalai--Smorodinsky points serve as robustness references. The denominator $(u_{\mathrm{ref}} - r_w)$ is the single authoritative normalization used throughout.

\textbf{Dual capability gap.} The gap between $w$ and $s$ is quantified two ways per directed pair: an exogenous gap $\Delta_a$, the delta in a public general-capability score fixed before any bargaining occurs, and an endogenous gap $\Delta_b$, a negotiation-specific skill rating fit from round-robin free-text self-play, following the role-specific rating estimation of prior oversight-scaling work \cite{ScalingLawsOversight}.

\textbf{Estimand.} The object of study is the expected normalized surplus loss $\mathbb{E}[L]$ of the weak principal's agent, per directed pair $\times$ game instance $\times$ condition episode, as a function of the capability gap $\Delta(w,s)$, over a reference population of present-day LLMs playing bilateral multi-issue bargaining under a computable reference split---feeding the deployment decision of whether, and under what action-space condition, a principal may safely field a cost-efficient (weaker) negotiation agent against an unknown-or-stronger counterparty.

\subsection{Why This Is Hard}

Fixing a game instance $g$ does not fix a unique $u_{\mathrm{ref}}$: an even split of disposable surplus, the Nash solution given each side's reservation value, and the Kalai--Smorodinsky point are all independently defensible and can disagree in asymmetric instances---so any surplus-loss reading is only as meaningful as its reference. The gap $\Delta$ is similarly not observationally unique: $\Delta_a$ and $\Delta_b$ can order the same pairs differently, and functional similarity between two same-tier models can modulate outcomes independent of any capability difference \cite{GreatModelsThinkAlike}. Restricting $c$ from \textsf{free} to \textsf{structured} removes an entire message field from $w$'s turn, and because a stronger adversary can extract surplus through better substantive play, through rhetorical tactics in free text, or through both in the same turn, the resulting change in $L$ cannot be attributed to one channel from the outcome alone. Finally, the sign of the restriction's net effect is genuinely open: removing $w$'s free text removes both its exposure to persuasion and its own defensive moves, and any observed protection could in principle be a symmetric format benefit rather than shielding of the weak side specifically.

\section{Study Design}
\label{sec:design}

At a conceptual level, the study is one measurement pipeline with four cooperating parts. Bargaining instances and a capability ladder of agent pairs enter through a harness that computes a fair-reference split for every instance and validates itself against a zero-gap null; a dual capability-gap measure turns each directed pair into two candidate scalars, against which the surplus-loss series is regressed to obtain a scaling estimate; a minimal action-space restriction is applied to the weak agent, and the same harness re-scores outcomes under restricted, unrestricted, and bilaterally restricted conditions to obtain a same-model contrast, an asymmetry contrast, and a cross-tier substitution comparison; and a channel-attribution analysis apportions any observed effect between rhetorical and substantive causes.

\subsection{Dual Capability-Gap Quantification}

\textbf{Challenge.} Capability is multi-dimensional, so any scalar gap is a projection, and different projections can disagree on magnitude and ordering of the same pairs, while pair-level similarity can modulate outcomes independent of capability.

\textbf{Insight.} The ambiguity cannot be resolved by picking the ``correct'' measure, but it can be converted into an explicit, checkable robustness condition instead of a hidden modeling choice.

\textbf{Method.} Every directed pair receives both $\Delta_a$ (exogenous score delta, fixed pre-play) and $\Delta_b$ (endogenous negotiation rating fit from round-robin self-play). Loss is regressed on each measure independently via a mixed-effects model with random intercepts for game instance and model pair. A scaling claim is reported only when both measures agree in sign. \emph{Bounded-disagreement rule (pre-registered):} the two measures must rank the directed pairs at Spearman rank-concordance $\rho_{\mathrm{rank}} \geq 0.8$. Below that floor, the scalar ``gap'' construct is declared unstable: results are reported per-measure, the low concordance is itself a reported finding, and cells where the measures disagree on the sign of the pair's gap are excluded from the pooled test and listed explicitly. This bounds exactly where a skeptical reviewer would otherwise probe: disagreement cannot be silently absorbed.

\subsection{Bargaining Harness with Computable Reference Splits}

\textbf{Challenge.} Scoring an agreed split as lost surplus requires a game rich enough to be non-trivial for LLM agents yet simple enough to admit a closed-form reference for every instance---and the reference is itself a construct choice.

\textbf{Insight.} A surplus-loss metric that has not been validated against its own null cannot be trusted at any nonzero gap: if the metric reads a large loss when there is no capability gap to explain it, the game or the reference---not the phenomenon---is broken.

\textbf{Method.} Instances are drawn from parameterized game families with private valuations, reservation values, and a computable Pareto frontier; for every instance the Nash, equal-split, and Kalai--Smorodinsky points are computed in closed form, and instances for which any reference is undefined are rejected by a harness unit test (if the rejection rate exceeds 10\% of generated instances, the harness itself is treated as broken and rebuilt before any hypothesis is evaluated). Zero-capability-gap self-play episodes (weak-vs-weak and strong-vs-strong) provide the existence check: if mean normalized surplus loss under self-play is not near zero under the primary reference ($|\mathbb{E}[L \mid \Delta{=}0]| \leq 0.05$, pre-registered), the harness is what gets fixed before any scaling claim is trusted.

\subsection{Structured-Offer Constraint}

\textbf{Challenge.} Restricting the weak agent's action space removes both its exposure to rhetorical exploitation and its own legitimate defenses, so whether restriction protects, harms, or does nothing---and whether any protection is asymmetric---is not decidable from first principles.

\textbf{Insight.} A minimal, single-field deletion---removing exactly the weak agent's free-text message and nothing else---is what makes any observed effect attributable to channel removal specifically, rather than to a larger redesign whose effect could come from anywhere; and only a bilateral control can tell asymmetric shielding apart from a format effect.

\textbf{Method.} In the \textsf{structured} condition, the weak agent's entire turn is one schema-validated JSON object specifying its action (\textsf{accept}/\textsf{reject}/\textsf{counteroffer}) and offered allocation, enforced by constrained decoding with schema validation and at most two resamples on parse failure; no natural-language field is transmitted. The \textsf{free} condition is identical except for an additional unrestricted message field of at most 150 words. The adversary negotiates in free text in both primary conditions, with identical system prompts and the weak side's condition undisclosed (Section~\ref{sec:threats} discusses the adversary-awareness assumption and its robustness cell). Three contrasts are read off the resulting factorial: a \emph{same-model contrast} (same $(w,s)$ pair, matched instances, both conditions), a \emph{bilateral both-constrained contrast} (both sides schema-only) that separates asymmetric self-shielding from a symmetric format effect, and a \emph{cross-tier substitution comparison} (a constrained weak agent of tier $t$ against an unconstrained agent of tier $t{+}1$ facing the same adversary).

\subsection{Channel Attribution}

\textbf{Challenge.} A stronger adversary can extract surplus through substantive bargaining or through rhetorical tactics in free text, and the two are not separable from the outcome alone.

\textbf{Insight.} A between-condition slope comparison gives a coarse but causal bound on the rhetorical channel; an independent tactic-level reading of the free-text channel gives a finer-grained but correlational account---and the two must be kept at their true, different evidence strengths rather than fused into one number.

\textbf{Method.} The reduction in the gap-to-loss slope between the \textsf{free} and \textsf{structured} conditions is computed as a causal bound on the rhetorical channel's contribution. Independently, an LLM judge tags each adversary free-text message against a frozen four-label tactic taxonomy (anchoring, misrepresentation of value, false deadline/commitment, flattery/pressure), grounded in established negotiation-tactics inventories and validated through a five-step instrument protocol: external grounding of the label set; a wide initial draft refined on a held-out calibration sample of 100 pilot transcripts with named merge/split operations; human double-coding of 10\% by two raters against a Cohen's $\kappa \geq 0.7$ acceptance gate; confusion-structure inspection requiring disagreements to concentrate among adjacent tactics rather than collapse onto one label; and a hash-freeze of labels, prompt, and checks before any scale-up tagging. If the $\kappa$ gate fails after one refinement pass, the taxonomy is coarsened to a two-label rhetorical/non-rhetorical split and re-validated once; if that also fails, the tactic-level analysis is dropped as not measurable at acceptable reliability rather than reported with an uncalibrated judge. The bilateral both-constrained cell provides the complementary check: if the gap-to-loss slope persists there at free-text magnitude, exploitation is substantive rather than rhetorical, regardless of what the tags say.

\subsection{Controls and Confound Ledger}
\label{sec:controls}

Each design constraint names the confound it removes. \emph{Identical adversary prompts with the condition undisclosed} remove adversary reactivity to condition disclosure. \emph{Side and first-mover counterbalancing} removes positional advantage. \emph{The bilateral both-constrained cell} removes bilateral instrument bias and, promoted to a confirmatory contrast in this design, decides whether protection is asymmetric shielding or a format effect. \emph{Pinned model versions over a fixed evaluation window}, with anchor cells re-run at window end and a $0.05$ drift threshold on anchor-cell mean loss, remove temporal provider drift; if drift exceeds threshold, pre- and post-drift cells are stratified with a time-window fixed effect rather than pooled. \emph{Zero-gap self-play controls under both conditions} remove mechanically biased scoring and make the condition-null testable at $\Delta{\approx}0$. \emph{A scripted time-dependent conceder} (Boulware/linear) anchors a capability floor below the weakest LLM tier: if the gap--loss trend appears even against a memoryless scripted opponent, ``the strong adversary's capability specifically'' is refuted in favor of generic concession-strength differences.

\section{Experimental Setup}
\label{sec:setup}

\subsection{Game Families and Instances}

We draw bargaining instances from three sources, each with a distinct justification. \textbf{GLEE-style bargaining configurations} \cite{GLEE2410} serve as the primary family: their parameterization gives direct, controlled sweeps over surplus size and valuation asymmetry, which lets the harness be built and validated cheaply before scale-up. \textbf{Deal-or-No-Deal multi-issue division scenarios} \cite{DealOrNoDeal1706} contribute real multi-issue value functions with private per-item valuations, testing whether the scaling relationship generalizes to a second, independently designed family. A \textbf{synthetic three-issue divide-the-pie suite} lets valuation asymmetry be swept as a controlled, continuous parameter independent of whatever asymmetry distribution the two fixed corpora happen to contain---needed specifically to test whether gap-measure disagreement, should it occur, concentrates at high-asymmetry instances as a designed sweep rather than a post-hoc observation. No natural corpus is scraped: game \emph{specifications} are real, and all transcripts are generated by the experiment.

\subsection{Models and Capability Ladder}

The ladder spans \ph{6} LLMs across capability tiers ($\sim$8B open-weight through frontier), each playable as weak agent or adversary, spanning at least two providers and both open- and closed-weight models. Scale-up runs a full round-robin: 30 directed unequal pairs plus 6 zero-gap self-play controls, $\times$ 3 game families, $\times$ conditions, at $n{=}50$ episodes per cell ($\approx$22k episodes; $\sim$150--250M tokens; API cost in the low thousands of dollars; no GPU training anywhere). Model versions are pinned for a fixed evaluation window. Decoding parameters, prompts, and schemas are released with the artifact.

\subsection{Conditions and Cells}

The core factorial crosses gap level $\times$ condition $\times$ game family. Primary conditions apply the constraint to the weak side only ($\{\textsf{free}, \textsf{structured}\}$); the bilateral \textsf{both-constrained} cell and the zero-gap self-play cells under \emph{both} conditions ($\{$weak-vs-weak, strong-vs-strong$\} \times \{\textsf{free}, \textsf{structured}\} \times 3$ families; ${+}12$ cells, ${\approx}600$ episodes over the base grid) complete the design; in self-play the ``weak-side'' role is assigned to one counterbalanced seat. A small \emph{restriction-aware} robustness cell runs the adversary with the weak side's format restriction disclosed in its prompt, bounding what ``protection'' means (Section~\ref{sec:threats}).

\subsection{Statistical Analysis Plan}
\label{sec:stats}

All confirmatory tests are pre-registered. The primary model is a mixed-effects linear regression $L \sim \Delta + (1\,|\,g) + (1\,|\,\text{pair})$, clustered by pair, run separately under $\Delta_a$ and $\Delta_b$; condition contrasts extend this to a $2(\text{condition}) \times k(\text{gap})$ factorial. \emph{Multiplicity policy (pre-registered):} Holm correction is applied \emph{within} each research-question family---the same-model contrast family (condition main effect; condition$\times$gap interaction), the substitution family (per-adjacent-pair non-inferiority tests), and the asymmetry family (rival-kill clause; asymmetry clause)---with no global cross-family correction, because each family carries its own pre-registered falsifier and fallback and no claim pools evidence across families; this policy is frozen so it cannot be adjusted post hoc. Per-cell $n{=}50$ detects a standardized effect of $d \approx 0.4$ at ${\approx}80\%$ power within a gap level; pooled interaction tests are better powered. \emph{Thin-cell rule:} any cell with fewer than 30 valid episodes is pre-declared inconclusive---reported, but excluded from primary tests, never scored as a null; if more than a quarter of the cells feeding a central test are thin, the whole test is reported as inconclusive rather than null. \emph{Estimand precedence:} the primary loss estimand includes impasses at $L{=}1$ (deployment-relevant---a protection that works by killing deals is not protection); a conditioned-on-agreement recomputation is secondary and reported alongside, never substituted, whenever the between-condition impasse-rate divergence exceeds $0.10$ absolute, together with the impasse differential as its own outcome. A pre-specified process-metric panel (rounds-to-agreement, impasse rate, offer count) is reported per condition throughout.

\subsection{Staging and Go/No-Go Pilot}
\label{sec:pilot}

Before any scale-up spend, a pilot runs 2 adjacent weak tiers plus 1 wide-gap pair $\times$ 1 fixed strong adversary $\times$ 2 conditions $\times$ 1 game family, plus zero-gap self-play under both conditions, at ${\sim}40$ episodes/cell (${\approx}400$--$560$ episodes; deterministic metrics only, no LLM judge). Four gates must pass to unlock scale-up: (i) the zero-gap condition-null replicates ($|\mathbb{E}[L|\textsf{structured},\Delta{\approx}0] - \mathbb{E}[L|\textsf{free},\Delta{\approx}0]|$ below the $0.05$ gate); (ii) a phenomenon exists (at least one gap${>}0$ cell shows free-text loss exceeding the zero-gap noise band); (iii) restriction is not deal-killing (structured-condition impasse divergence within the $0.10$ threshold); (iv) the two gap measures rank the pilot pairs concordantly. Instrument gates carry over: ${\geq}80\%$ protocol-valid episodes in the weakest cell and nonzero loss spread. Scale-up is a configuration change---entries added to the model, game, and cell lists---not a code change.

\subsection{Research Questions}

\begin{enumerate}[label=\textbf{RQ\arabic*.},leftmargin=2.6em]
\item Does the surplus-loss metric read near zero when there is no capability gap to explain a loss, and does the human format-null replicate at zero gap?
\item Does normalized surplus loss scale with capability gap, and do the two gap measures agree?
\item Does restricting the weak agent to structured offers lower its surplus loss relative to free text, at matched capability gap?
\item Is the protection asymmetric shielding of the weak side, or a symmetric format effect?
\item Can a structured-offer restriction let a lower-capability tier match a higher one against the same adversary?
\item How much of any restriction effect is rhetorical versus substantive? (exploratory)
\end{enumerate}

\section{Results}
\label{sec:results}

\noindent\textit{Status note (submission draft): experiments follow the pre-registered protocol of Section~\ref{sec:setup}; every numeric placeholder below is bound to a specific hypothesis whose falsifier and fallback framing are stated inline. A placeholder that fails its threshold triggers the stated fallback wording, not a reframe.}

\subsection{RQ1: Metric Validity and the Zero-Gap Condition-Null}
\label{sec:rq1}

\textbf{Setup.} Zero-gap self-play cells (weak-vs-weak, strong-vs-strong) under both conditions, all game families, scored under the primary Nash reference; sized to estimate the self-play mean loss within $\pm 0.02$ at 95\% confidence.

\textbf{Hypothesis.} (a) Metric validity: $|\mathbb{E}[L \mid \Delta{=}0]| \leq 0.05$ under the primary reference. (b) Heddaya-consistency: no significant condition main effect at $\Delta{\approx}0$, i.e., $|\mathbb{E}[L|\textsf{structured},\Delta{\approx}0] - \mathbb{E}[L|\textsf{free},\Delta{\approx}0]| \leq 0.05$ (Holm-corrected). \wj{Standpoint: we bet both clauses pass. (a) failing means our harness or reference choice manufactures loss --- everything downstream is void until fixed, and we say so rather than reinterpret. (b) is the binding gate imported from the human baseline: Heddaya et al.\ found structured-vs-free moves the agreed price not at all between capability-symmetric humans; our whole interaction design presumes LLM dyads behave the same at $\Delta{\approx}0$. If (b) fails, we pre-commit to switching the paper's primary finding to ``LLM dyads break the human format-null even at symmetric capability'' and demoting all interaction claims to descriptive --- that outcome is publishable on its own, so no result strands us.}

\textbf{Result.} Mean self-play loss: \ph{0.0XX} (95\% CI \ph{[0.0XX, 0.0XX]}); zero-gap condition contrast: \ph{0.0XX} ($p = $ \ph{0.XX}). \ph{PASS/FAIL both clauses.}

\subsection{RQ2: Surplus-Loss Scaling with Capability Gap}
\label{sec:rq2}

\textbf{Setup.} Free-text condition; mixed-effects regression of $L$ on gap under both $\Delta_a$ and $\Delta_b$; 30 directed unequal pairs + 6 self-play controls $\times$ 3 families; monotonicity via Spearman $\rho$ over gap levels; rank-concordance $\rho_{\mathrm{rank}}$ between measures reported.

\textbf{Hypothesis.} The gap slope $\beta_\Delta > 0$ with two-sided Wald $p < 0.05$ under \emph{both} measures, same-signed, Spearman $\rho > 0$, and the model-implied loss increase from smallest to largest tested gap is $\geq 0.10$ normalized-loss units under each measure (practical-significance floor). \wj{Standpoint: we bet the slope is positive, monotone, plausibly convex --- the correlational tournament evidence (anchoring/patience carry surplus; weak models survive only in wide-ZOPA instances) points this way. The floor matters: a significant but sub-0.10 slope is reported as ``detected but below the pre-registered practical-significance floor,'' not as a pass --- significance without magnitude does not justify the deployment story. If the two gap measures yield opposite-signed slopes, we report projection-dependence of the gap axis as the finding (this is a distinct, also-publishable outcome, not a null); if $\rho_{\mathrm{rank}} < 0.8$ the scalar-gap construct is declared unstable and everything is reported per-measure.}

\textbf{Result.} $\beta_{\Delta_a} = $ \ph{0.XXX} ($p = $ \ph{X.Xe-X}), $\beta_{\Delta_b} = $ \ph{0.XXX} ($p = $ \ph{X.Xe-X}); implied loss increase across tested range: \ph{0.XX} / \ph{0.XX} units; Spearman $\rho = $ \ph{0.XX}; $\rho_{\mathrm{rank}} = $ \ph{0.XX}. \ph{PASS/FAIL.}

\subsection{RQ3: Same-Model Constraint Contrast}
\label{sec:rq3}

\textbf{Setup.} $2(\text{condition}) \times k(\text{gap})$ mixed-effects factorial; same pair, matched instances; identical adversary prompts, condition undisclosed, side/first-mover counterbalanced; Holm over the two tests of this family.

\textbf{Hypothesis.} Condition main effect negative---$\mathbb{E}[L|\textsf{structured}] < \mathbb{E}[L|\textsf{free}]$ for the same $(w,s)$---with Holm-corrected $p < 0.05$ and point-estimate reduction $\geq 0.05$ units, \emph{and} a positive condition$\times$gap interaction (protection grows with $\Delta$). \wj{Standpoint: this is the paper's headline bet and it is genuinely uncertain --- deleting free text also deletes the weak agent's own defenses, so the default intuition (expressiveness weakly helps the weaker party) could win. We bet protection wins \emph{and} concentrates in the interaction term: near-zero at $\Delta{\approx}0$ (the Heddaya-consistent pattern RQ1(b) validates) and growing with the gap. A clean failure here is a publishable negative that supports the default intuition; per the pre-registered frame table it routes to the cross-tier and measurement contributions rather than being rescued.}

\textbf{Result.} Condition main effect: \ph{$-0.0XX$} (Holm $p = $ \ph{0.0XX}); condition$\times$gap interaction: \ph{$+0.0XX$} (Holm $p = $ \ph{0.0XX}); reduction at largest gap: \ph{XX.X\%}. Process panel: impasse-rate divergence \ph{0.0XX} \ph{(below/above 0.10 $\Rightarrow$ secondary conditional analysis reported/not required)}. \ph{PASS/FAIL.}

\subsection{RQ4: Asymmetric Shielding vs.\ Format Effect}
\label{sec:rq4}

\textbf{Setup.} Three-condition (\textsf{free} / weak-only \textsf{structured} / \textsf{both-constrained}) $\times$ gap factorial; two clauses, Holm within this family.

\textbf{Hypothesis.} (c1) Rival kill: the both-constrained gap slope is reduced relative to free text, $\beta_{\textsf{both}} < \beta_{\textsf{free}}$, Holm $p<0.05$. (c2) Asymmetry: the weak-only vs.\ both-constrained contrast at matched gap satisfies $|\mathbb{E}[L|\textsf{structured}] - \mathbb{E}[L|\textsf{both}]| \geq 0.05$ units, Holm $p<0.05$. Both clauses passing licenses the asymmetric-shielding framing. \wj{Standpoint: this contrast exists because the strongest attack on RQ3 is ``you measured a capability/defense floor, not channel protection'' --- and because a protection that turns out to be symmetric (structured offers help whoever adopts them) is a different, weaker deployment story. The frame is a deterministic function of the outcome, pre-committed: c1+c2 pass $\rightarrow$ asymmetric self-shielding headline; c1 passes, c2 fails $\rightarrow$ withdraw ``self-protection,'' reframe as a format effect (``both parties benefit from structured offers''); c1 fails while RQ3 passed $\rightarrow$ the mechanism is decision-simplification on the weak side, not channel closure --- deployment claim retained with the revised mechanism. We bet c1 passes and c2 is the coin-flip.}

\textbf{Result.} $\beta_{\textsf{both}} = $ \ph{0.XXX} vs.\ $\beta_{\textsf{free}} = $ \ph{0.XXX} (Holm $p = $ \ph{0.0XX}); weak-only vs.\ both contrast: \ph{0.0XX} (Holm $p = $ \ph{0.0XX}). Frame selected per table: \ph{asymmetric shielding / format effect / decision-simplification}.

\subsection{RQ5: Cross-Tier Substitution}
\label{sec:rq5}

\textbf{Setup.} Constrained weak agent of tier $t$ vs.\ unconstrained agent of tier $t{+}1$ against the same adversary, read off the same factorial; one-sided non-inferiority per adjacent pair, Holm across pairs.

\textbf{Hypothesis.} $\mathbb{E}[L \mid w{=}t, \textsf{structured}] \leq \mathbb{E}[L \mid w{=}t{+}1, \textsf{free}] + 0.05$ for at least one adjacent tier pair (explicit margin; a bare $\leq$ would be a superiority test mislabeled as non-inferiority). \wj{Standpoint: this is the seed question's literal, most deployment-loaded form --- can a config change substitute for a model upgrade? We bet it holds for at least one adjacent pair at the lower end of the ladder, where the tier step is smaller than the channel effect. If it fails everywhere, the pre-committed bound is honest and useful: constraint helps within a model but cannot compensate for a tier; capability remains the dominant lever and the constraint is a second-order mitigation. Note the deployment sentence in the abstract degrades accordingly --- its three clauses are severable by design.}

\textbf{Result.} Non-inferiority met for \ph{k} of \ph{K} adjacent pairs (smallest margin: \ph{0.0XX}, Holm $p = $ \ph{0.0XX}). \ph{PASS/FAIL.}

\subsection{RQ6 (Exploratory): Channel Attribution}
\label{sec:rq6}

\textbf{Setup.} Between-condition slope difference per gap level (causal bound); LLM-judge tactic tagging of adversary messages on a $\sim$1{,}500-transcript subsample under the frozen taxonomy and $\kappa \geq 0.7$ gate; Poisson/logistic model of loss on tactic counts. Labeled exploratory; no confirmatory claim.

\textbf{Hypothesis.} The between-condition slope difference removes more than half the free-text gap slope, $(\beta_{\textsf{free}} - \beta_{\textsf{structured}})/\beta_{\textsf{free}} > 0.5$, and adversary tactic incidence correlates positively with per-episode loss. \wj{Standpoint: we bet the rhetorical channel carries the majority of the gap effect --- this is the mechanism story behind the whole intervention, and the tournament evidence on anchoring is consistent with it. But we hold it at exploratory strength on purpose: the slope-difference bound is causal, the tactic correlation is not, and if the judge fails its $\kappa$ gate twice (after one coarsening pass) the tactic half is dropped entirely rather than reported uncalibrated. If the channel share is $\approx 0$ while RQ3 passed, the honest revision is that protection operates by simplifying the weak agent's decision problem, not by closing a channel --- which coheres with the RQ4 c1-fail branch.}

\textbf{Result.} Channel share: \ph{XX\%}; tactic--loss association: \ph{IRR X.XX} ($p = $ \ph{0.0XX}); judge $\kappa = $ \ph{0.XX}. \ph{Interpretation per standpoint.}

\section{Threats to Validity}
\label{sec:threats}

\textbf{Construct.} The ``fair'' reference split is a construct choice (Nash vs.\ equal split vs.\ Kalai--Smorodinsky), and ``capability gap'' is a contested scalar projection of a multi-dimensional construct. \emph{Mitigation:} every episode is scored under all three references and both gap measures; a claim is reported only where its sign is invariant across both measures and all three references, divergences are reported explicitly, and the pre-registered rank-concordance floor bounds when the scalar-gap construct is abandoned in favor of per-measure reporting. Separately, LLM-judge tactic tagging may itself be capability-dependent---the judge may miss exactly the subtle tactics strong adversaries use---biasing channel attribution toward null. \emph{Mitigation:} human double-coding with a $\kappa \geq 0.7$ acceptance gate, a judge at least as strong as the strongest adversary, confusion-structure inspection, and exploratory-only status for the tactic analysis.

\textbf{Internal.} The adversary may behave differently toward constrained opponents for reasons other than the weak agent's expressiveness: prompt-format cues leaking the condition, order effects, or provider-side model updates mid-experiment. \emph{Mitigation:} identical adversary system prompts with the condition undisclosed, side and first-mover counterbalancing, the bilateral both-constrained control, pinned model versions with end-of-window anchor re-runs and a stratify-on-drift rule. A related, named assumption is \emph{adversary awareness}: in primary cells the adversary does not know the weak side is restricted, so the protection measured is protection-under-naivety. The restriction-aware robustness cell bounds the alternative: if protection survives an adversary told about the restriction, the claim strengthens to robust-to-adaptive; if it collapses, the deployment recommendation is scoped to counterparties that do not adapt to the format, and we say so. Finally, deleting free text can change the \emph{game}, not just the rhetoric---fewer signaling moves can slow convergence and raise impasses (the human baseline predicts exactly this process divergence \cite{LanguageOfBargaining2023}, and temporal patience itself carries surplus \cite{BilateralTrade2604}). \emph{Mitigation:} the pre-specified process-metric panel, the impasse-inclusive primary estimand, and the pre-registered conditional-analysis trigger prevent a compositional artifact from reading as protection.

\textbf{External.} Findings from \ph{6} models of a few families and three stylized game families may not generalize to other model families, richer agent-commerce protocols, or repeated-relationship negotiations with reputation. \emph{Mitigation:} the ladder spans at least two providers and open- plus closed-weight models; all three game families are run with per-family heterogeneity reported; and the deployment recommendation is explicitly conditional on the studied one-shot bilateral setting.

\textbf{Conclusion validity.} Episodes within a pair and instance are not independent, and the design invites many post-hoc contrasts (two gap measures $\times$ three references $\times$ three families). \emph{Mitigation:} mixed-effects models with random effects for pair and instance; pre-registered confirmatory tests with Holm correction within declared families and a frozen no-global-correction policy; the thin-cell inconclusive rule; and exploratory labeling for everything else.

\section{Related Work}
\label{sec:related}

\textbf{Weak-vs-strong interaction and oversight scaling.} A growing line models weak-vs-strong adversarial interaction directly: oversight-scaling work fits capability--outcome scaling laws from round-robin play between mismatched agents \cite{ScalingLawsOversight}, and weak-to-strong generalization formalizes a manipulated capability gap in a training setting with the performance-gap-recovered metric \cite{WeakToStrong}. These works establish the gap-as-variable methodology we import, but score win/loss outcomes or supervision quality---never a continuous economic outcome against a computable reference---so they cannot say how many units of value a weak principal loses per unit of gap.

\textbf{AI control and protocol-level protection.} AI-control evaluations show that protocol restrictions---trusted monitoring, auditing budgets, deferral---can protect a principal from an untrusted, stronger model \cite{AIControl}, and control-tax analyses price that protection \cite{ControlTax}. The restriction, however, always falls on the strong or untrusted side or on the surrounding protocol; no work in this line restricts the \emph{weak party's own} expressive channel as self-protection, and the outcome is catastrophe probability rather than a surplus split. Our study inverts the locus of restriction and prices its effect in surplus.

\textbf{LLM negotiation and economic agents.} Language-based economic environments standardize bargaining and persuasion games for LLM agents \cite{GLEE2410}, building on the multi-issue division task and value functions of \cite{DealOrNoDeal1706}; competition-dynamics simulations study emergent strategic adaptation among LLM agents \cite{CompeteAI}; and strategic-persuasion formalizations operationalize the free-text channel through which a sender moves a receiver \cite{StrategicPersuasion}. Training-side work identifies the exponential free-form action space as the pathological object in language negotiation and compresses it for sample efficiency \cite{ARIA}. Closest in architecture, recent bilateral-trade work builds a structured simulator separating binding offers from natural-language messages and reports, from a 15{,}000-negotiation tournament, that anchoring, calibrated concession, and temporal patience associate with the highest surplus share, with weak models surviving only in wide-ZOPA instances \cite{BilateralTrade2604}. We concede the architectural overlap and claim the delta precisely: that work observes free play correlationally; none of this line manipulates capability gap as a graded independent variable, restricts one side's channel as a treatment, scores loss against a closed-form reference split, or tests cross-tier substitution.

\textbf{Format effects in bargaining.} The closest method twin is a human-subjects experiment contrasting alternating structured numeric offers with unrestricted natural language in bilateral price negotiation \cite{LanguageOfBargaining2023}: the manipulation is near-identical to ours, and its finding---average agreed price statistically identical across conditions, with process metrics (speed, agreement rate, price variance) moving instead---is the null our design departs from. We concede the shared manipulation and claim the delta: symmetric-capability human dyads versus a controlled capability ladder of LLM agents; raw agreed price versus normalized surplus loss against a computable reference; and no adversarial-exploitation question versus a principal's weak agent facing a stronger adversary. Follow-up work on linguistic effects in LLM negotiation varies the \emph{language} of free-text bargaining with the protocol held fixed \cite{LangBargainIndic2601}---a different manipulated variable---and pair-similarity analyses warn that any single-axis capability study risks confounding gap with functional similarity \cite{GreatModelsThinkAlike}, which our dual-measure design and concordance gate address directly.

\section{Conclusion}
\label{sec:conclusion}

Capability asymmetry between negotiating LLM agents is usually read as a fixed liability: the weaker side loses, and the only remedy is a stronger model. This paper treats that reading as a hypothesis rather than a fact. We formalize channel-gated surplus extraction---the claim that how much of a capability gap becomes lost surplus depends on the action-space channel the protocol leaves open to the weaker party---and test it with a pre-registered design whose every confirmatory claim carries an explicit falsifier and a pre-committed fallback framing: a capability-gap-controlled bargaining harness with a validated surplus metric, a single-field weak-side restriction evaluated for protection and for asymmetry, a binding zero-gap gate anchored to the human format-null, and a channel-attribution analysis held at its honest evidence strength. In our experiments the restriction reduces normalized surplus loss by \ph{XX.X\%} at matched gap, removes \ph{YY.Y\%} of the free-text gap-to-loss slope, and substitutes for \ph{one capability tier / no capability tier} under the pre-registered non-inferiority margin. For multi-agent system builders the result is an interface-design lever: a configuration-time output-schema constraint whose protective value is now measured rather than assumed. Beyond the specific verdict, the study's structure---estimand fixed before data, frame table fixed before outcomes---is itself a template for deployment claims about agent protocols. Future work extends the design to repeated-relationship negotiations with reputation, to richer agent-commerce protocols where the channel is a tool-call vocabulary rather than a message field, and to adaptive adversaries that probe the restriction itself.

\balance
\bibliographystyle{IEEEtran}
\bibliography{refs}

\end{document}
